\section*{Likovi i njihove unutrašnje linije}

\subsection*{Konzul}

Motivacija: Želi da dovrši osvetu protiv Hegemonije zbog onoga što je učinila Maui-Kovenantu, njegovoj porodici i svetu čiji je gubitak od detinjstva nosio kao ličnu obavezu. Povratak na Hiperion za njega nije obična diplomatska misija već završetak puta koji je započeo zakletvom kod Sirine grobnice.\\\\

Strah: Ne plaši se samo Šrajka ni sopstvene smrti, već mogućnosti da je čitav njegov život bio deo tuđeg plana i da ni njegova izdaja ni osveta nisu bile istinski njegove. Plaši ga i saznanje da je, pokušavajući da kazni one koje smatra krivcima, možda pomogao da započne katastrofu mnogo veću od svega što je želeo.\\\\

Krivica / konflikt: Istovremeno je služio Hegemoniji koju je mrzeo, sarađivao sa Proteranima, izdao obe strane i na kraju sam aktivirao uređaj koji je počeo da uklanja antientropijska polja Vremenskih grobnica. Njegova osveta proistekla je iz stvarnih zločina Hegemonije, ali ga to ne oslobađa odgovornosti za sopstvene izbore. Posebno ga progone smrt supruge Greše i sina Alana, uništenje Maui-Kovenanta i saznanje da je godinama pomagao sistemu protiv kojeg se zakleo da će se boriti.\\\\

Promena kroz priču: Počinje kao zatvoren čovek koji krije sopstvenu ulogu i posmatra ostale hodočasnike sa diplomatske distance. Kada konačno ispriča istinu, prihvata mogućnost da ga grupa osudi ili ubije. Njihova odluka da ipak nastave zajedno prekida njegovu dugogodišnju izolaciju. Na kraju više ne razmišlja samo o osveti: kada vidi Rejčel, odlučuje da od Šrajka ipak nešto zatraži --- za nju.

\subsection*{Lenar Hojt}

Motivacija: U početku želi da otkrije sudbinu Pola Direa i izvrši dužnost koja mu je poverena. Posle susreta sa Bikurom njegova neposredna motivacija postaje mnogo ličnija: oslobađanje od kruciforme i bola koji ga prati čak i kroz fugu. Hodočašće Šrajku predstavlja mu poslednju mogućnost da pronađe izlaz iz stanja iz kojeg medicina i Crkva nisu uspele da ga izbave.

\subsection*{Fedman Kasad}

Motivacija: Želi ponovo da pronađe Monetu i Šrajka i da prekine vezu koja ga je godinama vodila između ljubavi, seksualne privlačnosti, nasilja i rata. Posle onoga što je doživeo na Hiperionu više ne dolazi da traži uslugu ili odgovor. Njegova namera je sasvim određena: ako ih ponovo sretne, pokušaće da ih ubije. Iza toga stoji i strah da su Moneta i Šrajk koristili upravo njega kao sredstvo za pokretanje većeg rata.

\subsection*{Martin Silenus}

Motivacija: Želi da dovrši \emph{Spev o Hiperionu}, jedino delo koje smatra istinski velikim i razlogom sopstvenog postojanja. Uveren je da je Šrajk njegova muza i da su se njegovi najbolji stihovi rađali dok je Šrajk ubijao ljude u Gradu Pesnika. Zato se vraća na Hiperion iako zna kakvu cenu taj povratak može imati. Njegova želja za umetničkom besmrtnošću neprestano je u sukobu sa spoznajom da je stvaranje njegove poeme možda povezano sa smrću drugih.

\subsection*{Sol Vejntraub}

Motivacija: Želi da spase Rejčel od Merlinove bolesti pre nego što se njen život povuče sve do rođenja i nestane. Hodočašće je njegov poslednji pokušaj nakon što su medicina, istraživanje i vera ostali bez odgovora. Ali njegova unutrašnja linija nije samo očinska žrtva: kroz snove o Abrahamu odbacuje ideju da Bog ima pravo da zahteva dete od roditelja. Na Hiperion ne vodi Rejčel da bi je poslušno žrtvovao, već da bi zahtevao odgovor i izbor od samog Boga.

\subsection*{Bron Lamija}

Motivacija: U početku želi da otkrije ko je pokušao da ubije Džonija i zašto nedostaje deo njegovog sećanja. Kako istraga postaje lična, njena motivacija se menja: želi da zaštiti Džonija, pomogne mu da postane ono što sam bira da bude i razotkrije sile iz TehnoCentra koje njime manipulišu. Posle njegove smrti nastavlja put na Hiperion noseći i njegovo nerođeno dete i Šrenovu petlju sa delom onoga što je Džoni bio. Njeno hodočašće zato postaje nastavak njegove nedovršene potrage za sopstvenim identitetom i razlogom zbog kojeg je bio povezan sa Hiperionom.

\subsection*{Het Mastin}

Uloga i misterija: Mastin je kapetan templarskog drvobroda Igdrasil, Pravi Glas Drveta i sedmi član hodočašća. Za razliku od ostalih, nikada ne stiže da ispriča svoju priču. Sa sobom nosi Mebijusovu kocku u kojoj se, po zaključku grupe, nalazi erg koji je verovatno nameravao da upotrebi pri susretu sa Šrajkom. Posle uništenja Igdrasila reaguje neobično uzdržano, zatim nestaje iz vetrokola, ostavljajući za sobom krv i kocku u neutralnom stanju. Kasnije ga hodočasnici iz Utvrđenja Hronos možda ponovo vide kako sam ide prema Vremenskim grobnicama. Prvi roman namerno ne objašnjava gde je bio, kako je prešao Masiv Uzde, šta je tačno nameravao sa ergom niti šta je želeo od Šrajka; njegova nedovršena priča ostaje jedna od otvorenih misterija na kraju knjige.